\
        %% (Volume 3) Chapter 5: Resonance and Meaning in Literary and Artistic Traditions

        \chapter{Resonance and Meaning in Literary and Artistic Traditions}
        \label{ch:v3c5}

        \epigraph{(Opening quotation from the existing literature on
        Idea Theory

EXISTING PROJECT SUMMARY: Existing project 'IdeaTheory' (5 volume(s), 11/62 tasks done). Volume titles: ['Mathematical Foundations of Idea Theory', 'Idea Theory and the Social Sciences', 'Idea Theory and the Humanities', 'Idea Theory, Cognitive Science, and Philosophy of Mind', 'Emergence: A Formal Theory']. Completed tasks: ['Lean toolchain and lakefile', 'Lean 4 foundations', 'Lean theorems (Vol 1): algebraic composition structure', 'Lean theorems (Vol 1): emergence and the cocycle condition', 'Lean theorems (Vol 1): resonance, weight, and Aufhebung decomposition', 'Lean theorems (Vol 1): asymmetry of resonance and meaning curvature', 'Lean theorems (Vol 1): idea structures, chains, and conjugation', 'Lean theorems (Vol 1): universal composition and foundational symmetry', 'Lean theorems (Vol 1): compositional identity and transmission chains', 'Lean theorems (Vol 1): advanced properties and multi-term reassociation', 'Lean theorems (Vol 1): idea structures and structural equivalence'].

STEERING INSTRUCTIONS: make all of theorems 1 through 9 just basic math, and then try to have 5 volumes aboutt reasonable applications to social sciences, humanities, philosophy of mind/cogsci, the 'problem' of emergence, and other advanced applications - the theory should still be just as impressive, but everything should have an interpretation in terms of these other areas

Regenerate the PRD to reflect the steering instructions. Preserve the pass/fail status of completed tasks where still applicable, but update, rename, or add tasks as directed. Rename/restructure volumes, lean files, website pages, and LaTeX chapters as needed.

EXISTING PROJECT SUMMARY: Existing project 'IdeaTheory' (6 volume(s), 11/54 tasks done). Volume titles: ['Mathematical Foundations of Idea Theory', 'Idea Theory and the Social Sciences', 'Idea Theory and the Humanities', 'Idea Theory, Cognitive Science, and Philosophy of Mind', 'Emergence: A Formal Theory', 'Advanced Applications and Unifying Perspectives']. Completed tasks: ['Lean toolchain and lakefile', 'Lean 4 foundations', 'Lean theorems: idea composition monoid', 'Lean theorems: resonance pairing', 'Lean theorems: Aufhebung decomposition', 'Lean theorems: emergence cocycle', 'Lean theorems: meaning curvature', 'Lean theorems: conjugation of ideas', 'Lean theorems: transmission chains', 'Lean theorems: structural equivalence', 'Lean theorems: graded idea algebra'].

STEERING INSTRUCTIONS: make all of theorems 1 through 9 just basic math, and then try to have 5 volumes aboutt reasonable applications to social sciences, humanities, philosophy of mind/cogsci, the 'problem' of emergence, and other advanced applications - the theory should still be just as impressive, but everything should have an interpretation in terms of these other areas

Regenerate the PRD to reflect the steering instructions. Preserve the pass/fail status of completed tasks where still applicable, but update, rename, or add tasks as directed. Rename/restructure volumes, lean files, website pages, and LaTeX chapters as needed.

EXISTING PROJECT SUMMARY: Existing project 'IdeaTheory' (6 volume(s), 11/54 tasks done). Volume titles: ['Mathematical Foundations of Idea Theory', 'Idea Theory and the Social Sciences', 'Idea Theory and the Humanities', 'Idea Theory, Cognitive Science, and Philosophy of Mind', 'Emergence: A Formal Theory', 'Advanced Applications and Unifying Perspectives']. Completed tasks: ['Lean toolchain and lakefile', 'Lean 4 foundations', 'Lean theorems: idea composition monoid', 'Lean theorems: resonance bilinear pairing', 'Lean theorems: Aufhebung decomposition', 'Lean theorems: emergence 2-cocycle', 'Lean theorems: meaning curvature form', 'Lean theorems: conjugation action', 'Lean theorems: transmission chain', 'Lean theorems: structural equivalence', 'Lean theorems: graded idea algebra'].

STEERING INSTRUCTIONS: make all of theorems 1 through 9 just basic math, and then try to have 5 volumes aboutt reasonable applications to social sciences, humanities, philosophy of mind/cogsci, the 'problem' of emergence, and other advanced applications - the theory should still be just as impressive, but everything should have an interpretation in terms of these other areas

Regenerate the PRD to reflect the steering instructions. Preserve the pass/fail status of completed tasks where still applicable, but update, rename, or add tasks as directed. Rename/restructure volumes, lean files, website pages, and LaTeX chapters as needed., relevant to Resonance and Meaning in Literary and Artistic Traditions.)}{--- (Author, Year)}

        \section{Background from the Literature}
        \label{sec:v3c5.lit}

        (300--600 word substantive narrative on what existing non-formalized
        writing on Idea Theory

EXISTING PROJECT SUMMARY: Existing project 'IdeaTheory' (5 volume(s), 11/62 tasks done). Volume titles: ['Mathematical Foundations of Idea Theory', 'Idea Theory and the Social Sciences', 'Idea Theory and the Humanities', 'Idea Theory, Cognitive Science, and Philosophy of Mind', 'Emergence: A Formal Theory']. Completed tasks: ['Lean toolchain and lakefile', 'Lean 4 foundations', 'Lean theorems (Vol 1): algebraic composition structure', 'Lean theorems (Vol 1): emergence and the cocycle condition', 'Lean theorems (Vol 1): resonance, weight, and Aufhebung decomposition', 'Lean theorems (Vol 1): asymmetry of resonance and meaning curvature', 'Lean theorems (Vol 1): idea structures, chains, and conjugation', 'Lean theorems (Vol 1): universal composition and foundational symmetry', 'Lean theorems (Vol 1): compositional identity and transmission chains', 'Lean theorems (Vol 1): advanced properties and multi-term reassociation', 'Lean theorems (Vol 1): idea structures and structural equivalence'].

STEERING INSTRUCTIONS: make all of theorems 1 through 9 just basic math, and then try to have 5 volumes aboutt reasonable applications to social sciences, humanities, philosophy of mind/cogsci, the 'problem' of emergence, and other advanced applications - the theory should still be just as impressive, but everything should have an interpretation in terms of these other areas

Regenerate the PRD to reflect the steering instructions. Preserve the pass/fail status of completed tasks where still applicable, but update, rename, or add tasks as directed. Rename/restructure volumes, lean files, website pages, and LaTeX chapters as needed.

EXISTING PROJECT SUMMARY: Existing project 'IdeaTheory' (6 volume(s), 11/54 tasks done). Volume titles: ['Mathematical Foundations of Idea Theory', 'Idea Theory and the Social Sciences', 'Idea Theory and the Humanities', 'Idea Theory, Cognitive Science, and Philosophy of Mind', 'Emergence: A Formal Theory', 'Advanced Applications and Unifying Perspectives']. Completed tasks: ['Lean toolchain and lakefile', 'Lean 4 foundations', 'Lean theorems: idea composition monoid', 'Lean theorems: resonance pairing', 'Lean theorems: Aufhebung decomposition', 'Lean theorems: emergence cocycle', 'Lean theorems: meaning curvature', 'Lean theorems: conjugation of ideas', 'Lean theorems: transmission chains', 'Lean theorems: structural equivalence', 'Lean theorems: graded idea algebra'].

STEERING INSTRUCTIONS: make all of theorems 1 through 9 just basic math, and then try to have 5 volumes aboutt reasonable applications to social sciences, humanities, philosophy of mind/cogsci, the 'problem' of emergence, and other advanced applications - the theory should still be just as impressive, but everything should have an interpretation in terms of these other areas

Regenerate the PRD to reflect the steering instructions. Preserve the pass/fail status of completed tasks where still applicable, but update, rename, or add tasks as directed. Rename/restructure volumes, lean files, website pages, and LaTeX chapters as needed.

EXISTING PROJECT SUMMARY: Existing project 'IdeaTheory' (6 volume(s), 11/54 tasks done). Volume titles: ['Mathematical Foundations of Idea Theory', 'Idea Theory and the Social Sciences', 'Idea Theory and the Humanities', 'Idea Theory, Cognitive Science, and Philosophy of Mind', 'Emergence: A Formal Theory', 'Advanced Applications and Unifying Perspectives']. Completed tasks: ['Lean toolchain and lakefile', 'Lean 4 foundations', 'Lean theorems: idea composition monoid', 'Lean theorems: resonance bilinear pairing', 'Lean theorems: Aufhebung decomposition', 'Lean theorems: emergence 2-cocycle', 'Lean theorems: meaning curvature form', 'Lean theorems: conjugation action', 'Lean theorems: transmission chain', 'Lean theorems: structural equivalence', 'Lean theorems: graded idea algebra'].

STEERING INSTRUCTIONS: make all of theorems 1 through 9 just basic math, and then try to have 5 volumes aboutt reasonable applications to social sciences, humanities, philosophy of mind/cogsci, the 'problem' of emergence, and other advanced applications - the theory should still be just as impressive, but everything should have an interpretation in terms of these other areas

Regenerate the PRD to reflect the steering instructions. Preserve the pass/fail status of completed tasks where still applicable, but update, rename, or add tasks as directed. Rename/restructure volumes, lean files, website pages, and LaTeX chapters as needed. says about \emph{Resonance and Meaning in Literary and Artistic Traditions}.  Identify the
        central informal claims, competing definitions, open questions, and
        folklore results.  Cite at least 4--8 real works, e.g.\ \cite{
        lean4_tpil} as a placeholder; replace with topic-appropriate
        citations from \texttt{references.bib}.)

        \section{From Informal Claims to Formal Statements}
        \label{sec:v3c5.bridge}

        Each informal claim above is rendered as a precise Lean 4 statement
        via the translation table below.  This is the connective tissue
        between the literature and the formalization.

        \begin{table}[ht]
        \centering
        \begin{tabular}{p{0.42\linewidth} p{0.28\linewidth} p{0.22\linewidth}}
        \toprule
        \textbf{Informal claim} & \textbf{Formal theorem} & \textbf{Source} \\
        \midrule
        (Informal claim 1) & \texttt{theorem\_v3c5\_1} & \cite{lean4_tpil} \\
        (Informal claim 2) & \texttt{theorem\_v3c5\_2} & \cite{mathlib4} \\
        (Informal claim 3) & \texttt{corollary\_v3c5\_1} & \cite{mac_lane_1998} \\
        \bottomrule
        \end{tabular}
        \caption{Translation table: informal $\rightarrow$ formal for
        \cref{ch:v3c5}.}
        \label{tab:v3c5.bridge}
        \end{table}

        \section{Definitions}

        \begin{definition}[Resonance and Meaning in Literary and Artistic Traditions]
        \label{def:v3c5.1}
          (Formal definition of the primary object of this chapter, with
          a citation to its informal antecedent in the literature.)
        \end{definition}

        \begin{remark}
          (Contextual remark: how does this definition compare to the
          informal definitions surveyed in
          \cref{sec:v3c5.lit}?)
        \end{remark}

        \section{Main Results}

        \begin{theorem}
        \label{thm:v3c5.1}
          (Statement of the main theorem of this chapter.)
        \end{theorem}

        \begin{proof}
          We proceed by the following steps.
          \begin{enumerate}
            \item (Step 1: establish the setup.)
            \item (Step 2: apply the key lemma.)
            \item (Step 3: conclude.)
          \end{enumerate}
          \qed
        \end{proof}

        \begin{corollary}
        \label{cor:v3c5.1}
          (Immediate corollary of \cref{thm:v3c5.1}.)
        \end{corollary}

        \begin{proof}
          Apply \cref{thm:v3c5.1} directly. \qed
        \end{proof}

        \section{Lean 4 Verification}

        \begin{lstlisting}[language=Lean4,caption={Lean proof of \cref{thm:v3c5.1}}]
        namespace IdeaTheory
        theorem theorem_v3c5_1 : True := by trivial
        end IdeaTheory
        \end{lstlisting}

        \section{Discussion: What the Formalization Reveals}
        \label{sec:v3c5.discuss}

        (Reflect on what the formal proofs above clarify, simplify, sharpen,
        or refute about the informal literature surveyed in
        \cref{sec:v3c5.lit}.  Cite further works.)

        \section{Notes and References}

        See \cite{lean4_tpil} for background on Lean 4 proof style.
